% Français 9115, batch 22 — Gallica views 421–440.
% Pass: Opus 5.5 (claude-opus-5-5), 2026-09-22 — first pass, unchecked against the views by a human.
%
% What the twenty views hold. One piece, begun before the batch and going
% on after it, in French, in her regular fair hand with few corrections
% (struck words, words added above the line, no marginal notes but one):
% a memoir of her own on the conditions at the limits of vibrating
% elastic bodies — first the lamina, in Euler's six cases, then the
% elastic plate — paginated by her (18) to (34) at the head of each page,
% with numbered paragraphs (N.o 8 at view 420, outside the batch; N.o 9 at
% view 433; « §. 4 », N.o 10 at view 439). Unlike the leaves of the
% critique in batches 11–15, these are written on both sides: the odd
% pages of her pagination carry the ink number at the top right, the even
% ones none. A small slip of calculations, mounted on a grey sheet, is
% bound in the middle (v. 428). No microfilm target, no printed matter.
%
%   v. 421  (18) the experiment she could not make directly; tuning forks
%           (« fourches métalliques ») obey the law of the lamina supported
%           at both ends although their ends are physically free (« V.
%           traité de Chladni N.o 88 P. 115 »); the conditions
%           dz/dx = 0, d³z/dx³ = 0 at such ends.
%   v. 422  (19) points satisfying d³z/dx³ = 0, dz/dx = 0 admitted among
%           the analytic limits, four more cases beside Euler's six, left
%           aside; Euler's equation « P. 120 » for free ends, z written
%           with u as variable, even and odd numbers of points of rest;
%           the proof announced that no point of rest can be a point of
%           support.
%   v. 423  (20) the proof: the couples z = 0, dz/dx; z = 0, ddz/dx² = 0
%           excluded (« suivant la remarque d'Euler N.o 47 »); yet the
%           middle of the lamina satisfies dz/dx = 0, d³z/dx³ = 0 and is an
%           analytic end.
%   v. 424  (21) conclusion for an even number of points of rest: two
%           half-laminae each free in the new sense at the middle; for an
%           odd number, the conditions at u = 1/2; the angle condition
%           (e^ω−1)/(e^ω+1) = tang ½ω.
%   v. 425  (22) that condition is Euler's second case (one end free, one
%           supported), which he need not have computed apart; the third
%           case (one fixed, one free) admits no inner limit; the fourth
%           (both supported), whose lamina bends like a stretched string.
%   v. 426  (23) there every point of rest is a true limit, 2n+1 limits
%           for n−1 points of rest; the fifth case contained in the sixth
%           (« v^eme », « vi^eme ») as the second in the first.
%   v. 427  (24) the analytic end always in the middle; points of rest of
%           laminae not to be likened to the nodes of strings; neither
%           Giordano « Ricati » (« Delle vibrazioni de cilindri T. 1 des M.
%           dela Société italienne ») nor « M^r Chladni » noticed that the
%           intervals between points of rest are unequal.
%   v. 428  a slip, not part of the fair copy: calculations without a
%           word, f², g², h², f²+g²−h², 4f²g², then a double integral of
%           (d²z/dx² + d²z/dy²) δ d²z/dx² dx dy integrated by parts. Ink
%           number 210.
%   v. 429  (25) what Ricati's values show (the part next to the end less
%           than half the next, the parts growing towards the middle);
%           Chladni would have seen it with longer laminae, « car je me
%           suis assuré que l'experience les confirme »; Euler did not see
%           it either — his N.o 47.
%   v. 430  the blank grey back of the mounting sheet of v. 428: not
%           transcribed.
%   v. 431  the leaf of v. 429 photographed a second time (same text, same
%           stamp, same « 211 »): not transcribed.
%   v. 432  (26) Euler's equation between the values of dz/dx on the two
%           branches holds for supported ends only; an obstacle that lets
%           the motion pass merely selects one mode, and for free or fixed
%           ends one does not hold the lamina at a point of support in
%           Euler's sense to get its sounds.
%   v. 433  (27) the vagueness of « libre » and « appuyé »; N.o 9, the
%           plates: from « la formule N.o 1 » the two limit equations S{…}
%           δ.dz/dx dy = 0, S d{…}/dx dy δz = 0; the rectangular plate.
%   v. 434  (28) the four couples of equations for the sides parallel to
%           the axis of y, the analogy with the lamina; z = 0 and dz/dx = 0
%           no longer mean a fixed point; the same for the sides parallel to
%           the axis of x.
%   v. 435  (29) the four couples for those sides; the corners, subject to
%           four conditions: « points de double limite ».
%   v. 436  (30) « points d'origine »; the round plate, x = r cos φ,
%           y = r sin φ, z = z' + z'', r = a at the limits.
%   v. 437  (31) the limit conditions in φ, satisfied by tang φ = 0 and
%           cot φ = 0, i.e. at the ends of the axes; elsewhere two
%           equations in z', z''.
%   v. 438  (32) the four couples that satisfy them; why the limit
%           conditions matter for the choice of particular integrals and
%           let one surface be conceived as several. The section ends two
%           thirds down the page.
%   v. 439  (33) « §. 4. Comparaison dela théorie à l'experience; dans le
%           cas des plaques quarrées », N.o 10: equation (C) of « N.o 3 »
%           with Euler's coefficients for free ends (« P. 118 et suiv. de
%           son m. »); straight nodal lines when one coefficient vanishes;
%           the supposition scriptA = −B and equation (C'); the first diagonal
%           nodal line x = y.
%   v. 440  (34) the second diagonal y = A − x for an even number of
%           points of rest, one only for an odd number; the nodal figure
%           passes through the intersections of the lines for scriptA = 0 and
%           B = 0. The page stops mid-sentence on « le nombre de ».
%
% Seams. View 421 opens a new paragraph on a new page, (18), after view
% 420 — page (17), ink 206, « N.o 8. Pour y parvenir… », which closes on
% « ce résultat. » — so nothing is cut at 420/421, but the piece began in
% batch 21. View 440 stops mid-sentence; view 441 (ink 216) continues it.
% References to pages outside the batch: « la formule N.o 1 » (v. 433),
% « l'équation (C) N.o 3 » (v. 439), « l'équation P. 120 dumémoire
% d'Euler » (v. 422).
%
% The numbers on the leaves, and why no \folio{} is written. (1) Her
% pagination, in parentheses at the head, centre: (18) at v. 421, then
% (19)–(24) at v. 422–427, (25) at v. 429, (26)–(34) at v. 432–440,
% without a gap. (25), (27), (29), (31), (33) are struck through with an
% oblique stroke; the others are not. (2) The ink series at the top
% right, in the larger rounder hand: 207 (v. 422), 208 (v. 424), 209
% (v. 426), 210 (v. 428, the slip), 211 (v. 429, repeated on v. 431), 212
% (v. 433), 213 (v. 435), 214 (v. 437), 215 (v. 439) — one per leaf,
% without a gap, on the side paginated odd, never on the even pages that
% are the backs of the same leaves; it continues 206 at v. 420 and goes on
% to 216 at v. 441. That it numbers leaves, not pages, and counts the
% small slip as a leaf, is what a foliation does; but as in every batch of
% this volume it is penned, not pencilled, and no \folio{} is written.
% (3) Her paragraph numbers: N.o 9 (v. 433), N.o 10 (v. 439), and « §. 4 »
% (v. 439). Every number above was read on its view; none is inferred.
% None of these leaves can be Laubenbacher and Pengelley's « Manuscript C »
% (ff. 348r–349r): the subject is elasticity, and the ink series here
% runs 207–215.
%
% Notation. Hers, kept. z the ordinate (often underlined in the prose, as
% are x, y, n, ω); ω the angle of Euler's cases; on v. 422–424 the
% variable of the lamina is written « u », kept as u. Her « S » is the
% sign of summation along the contour, kept as S; δ is the variation
% (her looped d). Round brackets round a partial derivative, braces and
% the points of suspension « = … » standing for an omitted factor are
% hers. Words she underlines in the prose are set in \emph{}; letters
% underlined inside the mathematics keep \underline{}. The first coefficient of v. 439–440 is a looped cursive capital,
% distinct from the A of the side of the square, and is set $\mathcal{A}$.
% « sin », « cos », « tang », « cot » are set \text{sin} etc.; the stop she
% writes after them is kept where it is clear. Where the leaf and the
% calculation disagree the leaf stands, with a note: v. 423 (a bracket
% (1−e^ω) where (1+e^ω) is carried down), v. 423 and 435 (a derivative
% written without « = 0 »), v. 428, v. 433, v. 437. Spelling is hers and
% is not flagged: « caractorise », « frapante », « parallelle »,
% « quarrée », « satisfesante », « Ricati », « cilindri », « apperçu »,
% « differens », « mouvemens », « dépendans », « résultans », « savans »,
% « raisonnemens », « longtems », « fesant », « séparés » for séparées,
% « remplis » for remplies, « dautant », and the run-together words
% « dela », « delalame », « aumilieu », « aunombre », « ceque », « desorte »,
% « adonné », « dumémoire », « parconséquent », « àprésent », « demême »,
% « deplus », « aulieu », « àpart ».
%
% What could not be read, and what is offered with doubt. \ill{} stands
% nine times: a sign blotted in an equation (v. 422); a struck word
% (v. 432); a struck group before « remarqué » (v. 433); the struck ending
% of « confirme » (v. 429); a sign after an integral and one after δ on
% the slip (v. 428); a struck word in the title of § 4 and a blackened
% group after « auroit » (v. 439); a struck word after « Il est »
% (v. 439). \uncertain{} stands for « qu'a » (v. 421, struck), « satisfont »
% and « le premier » (v. 421), « en » (v. 422), « caractorise » (v. 423),
% « remarque » (v. 425, struck), « d'extrêmités » (v. 426),
% « appartiendroient » (v. 434), « E. » and β (v. 433), the coefficient 2
% (v. 428), « r = a » (v. 436, struck), « fesant » (v. 439).
%
% Nothing here was checked against any published edition of Euler, of
% Chladni, of Riccati or of Sophie Germain's papers.
\documentclass[11pt,a4paper]{article}
% The preamble shared by every transcription of the mathematicians' manuscripts in Gallica.
%
% It rests on one idea: **uncertainty compounds, it does not resolve.** A
% transcription that smooths over a doubtful word has destroyed the only thing
% it was made for — a reader must be able to tell, at every point, what was on
% the page and what was supplied. The apparatus macros below are therefore the
% critical apparatus, not decoration.
%
% The same commands are recognised by `scripts/render.mjs`, which produces the
% site's reading view, and by `scripts/tei.mjs`. A macro added here must be
% added there too, or rendering fails loudly — which is the wanted behaviour.
%
% Comments here are English, like the rest of the repository. The documents
% this preamble sets are French, and that is deliberate: see the skills.

\ProvidesPackage{germain}[2026/09/15 Transcriptions of mathematicians' manuscripts in Gallica]

\usepackage{amsmath,amssymb,amsthm}
\usepackage{mathrsfs}
% Kept from the parent preamble so the renderer's diagram support stays
% available; these manuscripts draw no commutative diagrams, and nothing here uses it.
\usepackage{tikz-cd}
\usepackage[english,french]{babel}
\usepackage{fontspec}

% --- Greek ------------------------------------------------------------------
% The text face stays fontspec's default, Latin Modern, which has no Greek:
% XeTeX drops every glyph it lacks with a line in the log and nothing on the
% page, and the Greek words the manuscripts quote vanished from the PDFs. So
% the Greek and Greek Extended blocks get a XeTeX character class of their own,
% and entering or leaving a run of it switches the family to CMU Serif — the
% same Computer Modern design, with polytonic Greek — and back. Only the
% family changes: a Greek word in \textit or \textbf keeps its shape and series.
% The transitions to and from every other class carry the tokens that class
% has towards an ordinary letter, so French spacing before « ; » or inside
% « » still holds after a Greek word. No grouping, so no brace or \endgroup
% can come out unbalanced; maths is left alone, since XeTeX inserts nothing
% there. Kept identical in `exercices.sty`.
\makeatletter
\newfontfamily\germainGreekfont{cmun}[
  Extension=.otf, NFSSFamily=germaingreek,
  UprightFont=*rm, ItalicFont=*ti, SlantedFont=*sl,
  BoldFont=*bx, BoldItalicFont=*bi, BoldSlantedFont=*bl]
\newXeTeXintercharclass\germain@greek
\def\germain@greekrange#1#2{%
  \@tempcnta=#1\relax
  \loop
    \XeTeXcharclass\@tempcnta=\germain@greek
    \ifnum\@tempcnta<#2\advance\@tempcnta\@ne\repeat}
\germain@greekrange{"0370}{"03FF}
\germain@greekrange{"1F00}{"1FFF}
\def\germain@greekprev{\rmdefault}
\def\germain@greekin{%
  \edef\germain@greekprev{\f@family}\fontfamily{germaingreek}\selectfont}
\def\germain@greekout{\fontfamily{\germain@greekprev}\selectfont}
\def\germain@greekjoin{%
  \XeTeXinterchartokenstate=\@ne
  \@tempcnta=\z@
  \loop
    \ifnum\@tempcnta=\germain@greek\else
      \XeTeXinterchartoks\germain@greek\@tempcnta=\expandafter{%
        \expandafter\germain@greekout\the\XeTeXinterchartoks\z@\@tempcnta}%
      \XeTeXinterchartoks\@tempcnta\germain@greek=\expandafter{%
        \the\XeTeXinterchartoks\@tempcnta\z@\germain@greekin}%
    \fi
    \ifnum\@tempcnta<4095\advance\@tempcnta\@ne\repeat}
% babel-french sets its own classes, and saves and restores the interchar
% state, each time a language is selected: join again after it does.
\addto\extrasfrench{\germain@greekjoin}
\addto\extrasenglish{\germain@greekjoin}
\AtBeginDocument{\germain@greekjoin}
\makeatother

\usepackage{geometry}
\geometry{a4paper,margin=25mm,marginparwidth=28mm}
\usepackage{marginnote}
\usepackage{xcolor}
\usepackage[normalem]{ulem}
\setlength{\emergencystretch}{3em}
\usepackage{hyperref}
\hypersetup{colorlinks=true,linkcolor=black,urlcolor=black,pdfborder={0 0 0}}

% --- Batch metadata ---------------------------------------------------------
% Taken as they stand by the rendering script, which shows them at the head of
% the reading view. They fix what the file claims to cover. The macro names are
% the parent project's, so that the scripts stay shared; \folder here means the
% volume's slug — `fr-9115` — and \pages counts Gallica views.

\newcommand{\folder}[1]{\gdef\grothFolder{#1}}
\newcommand{\batch}[1]{\gdef\grothBatch{#1}}
\newcommand{\pages}[2]{\gdef\grothFirst{#1}\gdef\grothLast{#2}}
\newcommand{\foldertitle}[1]{\gdef\grothFoldertitle{#1}}

% The holder's own shelfmark, printed as they write it: « Français 9115 ».
\newcommand{\shelfmark}[1]{\gdef\grothShelfmark{#1}}

% The dating, copied verbatim from the holder's catalogue — never inferred
% here. « XIXe siècle » is the BnF's; a pass that dates a leaf from its content
% says so in a \note{}, as its own inference.
\newcommand{\dating}[1]{\gdef\grothDating{#1}}

% The status notice, printed in the title block and repeated as a diagonal
% watermark on every page of the PDF. The manuscripts are in the public
% domain; what this line declares is not a right but a quality — an
% unverified first pass by a machine. The skills write the line; a file
% without it gets no watermark, so leaving it out is visible in the output.
\newcommand{\watermark}[1]{\gdef\grothWatermark{#1}}

\gdef\grothFolder{?}\gdef\grothBatch{}\gdef\grothFirst{?}\gdef\grothLast{?}%
\gdef\grothFoldertitle{}\gdef\grothShelfmark{}\gdef\grothDating{}\gdef\grothWatermark{}

% --- The résumé -------------------------------------------------------------
\newenvironment{resume}
  {\small\setlength{\parskip}{0.45em}}
  {\par\medskip\hrule\medskip}

% --- Keywords ---------------------------------------------------------------
% In English, closing the modernised reading's résumé: the modern vocabulary
% under which to search for what the leaves construct. The manifest extracts
% them to tag the volume — this line is the single source of the tags.
\newcommand{\keywords}[1]{\par\smallskip\noindent
  {\footnotesize\textsc{Keywords} --- \textit{#1}}\par}

% --- The critical apparatus -------------------------------------------------

% The Gallica view number, in the margin. This is the anchor that lets a
% disagreement over a formula be settled by looking at *one* image rather than
% a volume of 750. The site uses it to turn the facsimile as the transcription
% is scrolled. A view is one scanned image: usually one side of a leaf.
\newcommand{\page}[1]{%
  \marginnote{\footnotesize\color{gray!70}\textsf{v.\,#1}}%
  \label{page:#1}\ignorespaces}

% The BnF's pencilled foliation, where the leaf carries it and a pass can read
% it: « 198r ». This is what the literature cites — « ff. 198r–208v » — and
% the one line that turns a citation into a view. Recorded, never inferred:
% a folio a pass cannot read is not written.
\newcommand{\folio}[1]{%
  \marginnote{\footnotesize\color{gray!50}\textsf{f.\,#1}}\ignorespaces}

% The modernised reading's coarser counterpart to \page{}: which views a
% section is drawn from.
\newcommand{\pagerange}[2]{%
  \marginnote{\footnotesize\color{gray!70}\textsf{v.\,#1--#2}}%
  \ignorespaces}

% Illegible: never guessed.
\newcommand{\ill}{\textcolor{red!60!black}{[\,\ldots\,]}}

% A doubtful reading: the word is given, and flagged as uncertain.
\newcommand{\uncertain}[1]{\uline{#1}}

% An editorial addition — a missing letter, an implied word.
\newcommand{\add}[1]{\textcolor{blue!50!black}{[#1]}}

% Struck out by the writer. What was crossed out is part of the document.
\newcommand{\struck}[1]{\sout{#1}}

% The transcriber's note, on the state of the page or the mathematics.
\newcommand{\note}[1]{\par\smallskip{\small\color{gray!60!black}\textit{[#1]}}\par\smallskip}

% A marginal note of hers, distinct from the body of the text.
\newcommand{\marginal}[1]{\par\smallskip\noindent\hspace{1em}%
  {\small\color{gray!60!black}#1}\par\smallskip}

% --- Title ------------------------------------------------------------------
% The title block is French, like the documents themselves.

\AddToHook{shipout/background}{%
  \ifx\grothWatermark\empty\else
    \begin{tikzpicture}[remember picture, overlay]
      \node[rotate=52, scale=3.4, align=center, text=gray!18,
            font=\sffamily\bfseries]
        at (current page.center) {\grothWatermark};
    \end{tikzpicture}%
  \fi}

\AtBeginDocument{%
  \begin{center}
    {\large\bfseries Manuscrits de mathématiciens (Gallica) ---
      \ifx\grothShelfmark\empty\grothFolder\else\grothShelfmark\fi}\\[2pt]
    {\itshape\grothFoldertitle}\\[4pt]
    {\small vues \grothFirst--\grothLast
      \ifx\grothBatch\empty\else\ (lot \grothBatch)\fi}\\[2pt]
    \ifx\grothDating\empty\else
      {\small\color{gray!60!black}Datation du catalogue : \grothDating}\\[2pt]
    \fi
    \ifx\grothWatermark\empty\else
      {\small\bfseries\color{red!45!gray}\grothWatermark}\\[2pt]
    \fi
    {\footnotesize\color{gray!60!black}Transcription établie d'après les images IIIF de Gallica.
     Source gallica.bnf.fr / Bibliothèque nationale de France.\\
     Les lectures incertaines sont soulignées, les illisibles notées [\,\ldots\,],
     les ajouts entre crochets ; \textsf{v.} numérote les vues de Gallica,
     \textsf{f.} la foliotation au crayon de la BnF.}
  \end{center}
  \vspace{1em}}

\endinput


\folder{fr-9115}
\shelfmark{Français 9115}
\batch{22}
\pages{421}{440}
\dating{1801-1900}
\watermark{Lecture automatique\\non vérifiée}
\foldertitle{Papiers de Sophie GERMAIN. — Recueil de dissertations et problèmes mathématiques et physiques.}

\begin{document}

\note{Numérotation des feuillets. Chaque page de cette mise au net porte
en tête, au milieu, entre parenthèses, sa pagination : (18) à (24) aux
vues 421 à 427, (25) à la vue 429, (26) à (34) aux vues 432 à 440, sans
lacune ; (25), (27), (29), (31) et (33) sont barrés d'un trait oblique.
Les feuillets sont écrits des deux côtés, et seul le côté de pagination
impaire porte en haut à droite, à l'encre, d'une écriture plus ronde, un
nombre de la série qui court dans tout le volume : 207 (vue 422), 208
(424), 209 (426), 211 (429, repris à la vue 431), 212 (433), 213 (435),
214 (437), 215 (439) ; le petit feuillet de calculs de la vue 428 porte
210. Cette série numérote des feuillets et non des pages, et elle compte
le petit feuillet ; elle occupe la place de la foliotation de la
Bibliothèque, mais elle est tracée à la plume et non au crayon, et aucun
« f. » n'est donc porté. Tous ces nombres ont été lus sur la vue ; aucun
n'est déduit. La vue 430 est le dos gris et blanc de la feuille qui porte
le feuillet de la vue 428 ; la vue 431 photographie une seconde fois le
feuillet de la vue 429 ; ni l'une ni l'autre n'est transcrite.}

\page{421}

On verra plus bas que j'ai été plus heureux par rapport a
un fait du même genre qui pourroit servir de preuve à celui
dont je n'ai pû m'assurer directement; et l'on ne s'étonnera
pas de la difficulté que j'ai trouvée à faire l'experience
dont il s'agit, si on fait attention \struck{\uncertain{qu'a}} a l'influence qu'ont
les moindres circonstances physiques pour exclure une partie des
mouvemens que l'analyse montre possibles et pour les
remplacer par d'autres \add{mouvemens} pris également aunombre de ceux
que l'analyse indique possibles, mais qui n'ont lieu
pourtant qu'à l'aide de la condition physique qui
n'entre pas dans le calcul. C'est ainsi que les fourches
\marginal{métalliques} qui doivent etre assujetties aux mêmes conditions analytiques
que les lames droites \add{, quoique la courbure ait pourtant un effet que j'expliquerai dans la suite,} donnent des sons qui suivent la
même loi que ceux qui appartiennent à la lame appuyée
a ses deux extremités, \struck{quoi} \add{malgré} que celles de la fourche soient
l'une et l'autre physiquement libres (V. traité de Chladni N.o 88
P. 115); desorte qu'il paroît que \add{les extremités} dela fourche \uncertain{satisfont}
veritablement au\add{x} couple\add{x} d'équations $\frac{dz}{dx}=0$, et $\frac{d^3z}{dx^3}=0$;
et que l'experience refuse de presenter pour elles le
cas $\frac{d^3z}{dx^3}=0$ et $\frac{d^2z}{dx^2}=0$; qui sur la lame droite
exclut presque toujours \uncertain{le premier}\struck{, par la facilité
que cette lame paroît avoir à se ployer aux
courbures qui appartiennent au second}.
\note{En tête, au milieu, « (18) », entre parenthèses, non barré ; aucun
nombre en haut à droite de cette vue. Écriture régulière et posée, la
main de ses mises au net, avec quelques corrections. « métalliques » est
écrit seul dans la marge de gauche, en face de « qui doivent etre
assujetties », sans signe de renvoi : il se rapporte apparemment à
« fourches ». L'addition « quoique la courbure ait pourtant un effet que
j'expliquerai dans la suite » est écrite en plus petit au-dessus de la
ligne, après « droites ». Le mot biffé devant « a l'influence » est
empâté ; lu « qu'a » avec doute. Le verbe après « fourche » est surchargé
à sa fin : on croit lire « satisfont » écrit sur « satisfait », ce
qu'appellent « les extremités » ajouté au-dessus et les « x » ajoutés à
« au » et à « couple » ; « au couple » devient ainsi « aux couples ».
Dans $\frac{d^3z}{dx^3}$ de la première équation, le $x$ du dénominateur
est tracé comme un « 1 ». « le premier » est écrit sur un autre mot,
peut-être « l'autre », et n'est pas sûr. Les deux dernières lignes et
demie, de « par la facilité » à « au second », sont barrées d'un trait
horizontal ; le point final est resté après « second ».}

\page{422}

Quoiqu'il en soit de la possibilité d'obtenir les figures et les
sons resultant des extremités assujetties aux conditions $\frac{d^3z}{dx^3}=0$ $\frac{dz}{dx}=0$,
il est impossible de ne pas admettre les points qui y satisfont
au nombre des limites analytiques des lames vibrantes; et rien
ne seroit plus aisé que d'obtenir, presque sans calcul, les valeurs
de $\underline{z}$ et de l'angle $\omega$ qui conviendroient aux quatre cas
que la considération de pareils points ajouteroit aux six
autres cas examinés par Euler. C'est à regret, et pour ne pas
m'écarter trop du sujet que j'ai a traiter, que j'abandonne
cette recherche; mais je ne puis me dispenser de considerer
aumoins les six cas dont Euler adonné l'analyse.

Je prends donc d'abord l'équation P. 120 dumémoire d'Euler.
En n'écrivant pas le coëfficient relatif au tems, \uncertain{en} changeant $y$ en
$z$, elle devient $z = \ldots\, e^{u\omega} \mp e^{(1-u)\omega} + (1\pm e^{\omega})\,\text{sin}.u\omega + (1\mp e^{\omega})\,\text{cos}.u\omega$.
Cette équation appartient au cas où les extrêmités satisfont
aux conditions $\frac{d^2z}{dx^2}=0$ et $\frac{d^3z}{dx^3}=0$; Elle présente elle même deux
cas, \struck{suivant} \add{selon} qu'il y a un nombre pair ou un nombre
impair de points de repos. Le premier se rapporte aux
signes inferieurs, et les signes superieurs conviennent au second.

Je \struck{vais commencer par prouver} \add{prouverai d'abord} que, lorsque l'on a
$z = \ldots\, e^{u\omega} + e^{(1-u)\omega} + (1\,\ill{}\,e^{\omega})\,\text{sin}\,u\omega + (1+e^{\omega})\,\text{cos}\,u\omega$, aucun point
de repos ne peut être consideré comme point d'appui
\note{En tête, au milieu, « (19) » ; en haut à droite, à l'encre, « 207 ».
La lettre soulignée est $z$ ; l'angle est $\omega$, tracé comme ailleurs
dans ce volume. La variable des exposants et des arcs est écrite « u » ;
elle est gardée telle quelle. « sin » est écrit avec un s long, suivi
d'un point à la première équation et non à la seconde. Le mot devant
« changeant » est empâté, peut-être corrigé : lu « en » avec doute.
L'initiale de « Elle présente » est repassée. Le signe du second
crochet de la seconde équation est effacé sous une tache d'encre ; les
deux autres signes y sont « + », là où la première équation porte
« $\mp$ » et « $\mp$ ». « suivant » est barré d'un trait épais,
« selon » écrit au-dessus. « vais commencer par prouver » est barré d'un
trait horizontal, « prouverai d'abord » écrit au-dessus. La phrase se
poursuit à la vue suivante.}

\page{423}

au sens d'Euler, ni même comme extrêmité analytique. En
effet, pour qu'un tel point satisfît aux conditions des
limites, il faudroit que l'un ou l'autre des couples d'équations
$z=0$, et $\frac{dz}{dx}$; $z=0$ et $\frac{ddz}{dx^2}=0$ lui convint. or le premier
\struck{indique} \uncertain{caractorise} le point fixe; et, suivant la remarque
d'Euler \struck{P. 157} \add{N.o 47} de son memoire, toute communication de
mouvement est impossible entre les deux portions separés
par un tel point. A l'égard du second couple, il se
reduit aux deux equations
\[ e^{u\omega} + e^{(1-u)\omega} + (1-e^{\omega})\,\text{sin}.u\omega + (1+e^{\omega})\,\text{cos}\,u\omega = 0 \]
\[ e^{u\omega} + e^{(1-u)\omega} - \{(1-e^{\omega})\,\text{sin}.u\omega + (1+e^{\omega})\,\text{cos}\,u\omega\} = 0 \]
desquelles on tire celles ci: $e^{u\omega} + e^{(1-u)\omega} = 0$, $(1-e^{\omega})\,\text{sin}\,u\omega + (1+e^{\omega})\,\text{cos}\,u\omega = 0$
dont la première est visiblement impossible et dont la seconde
ne peut s'accorder, dans aucun cas, avec l'équation qui,
dans l'analyse d'Euler, sert à déterminer l'angle $\omega$.

\struck{Cependant} \add{Néanmoins,} dans la manière de vibrer dont je m'occupe
en ce moment, un des points delalame satisfait aux
conditions des limites et peut parconséquent etre consideré
comme une extrêmité analytique. \struck{c'est le point de} Ce
point est situé au milieu delalame; car des conditions
$\frac{dz}{dx}=0$ et $\frac{d^3z}{dx^3}=0$ qui se réduisent aux deux équations
\[ e^{u\omega} - e^{(1-u)\omega} + (1-e^{\omega})\,\text{cos}\,u\omega + (1+e^{\omega})\,\text{sin}\,u\omega = 0 \]
\[ e^{u\omega} - e^{(1-u)\omega} - \{(1-e^{\omega})\,\text{cos}\,u\omega + (1+e^{\omega})\,\text{sin}\,u\omega\} = 0 \qquad \text{ou} \]
\[ e^{u\omega} - e^{(1-u)\omega} = 0,\quad (1-e^{\omega})\,\text{cos}\,u\omega + (1-e^{\omega})\,\text{sin}\,u\omega = 0 \]
l'une devient
identique par la supposition $u=\frac{1}{2}$ tandis que l'autre est
précisément l'équation qui sert à déterminer l'angle $\omega$.
\note{En tête, au milieu, « (20) » ; aucun nombre en haut à droite. La
phrase continue celle de la vue 422 (« aucun point de repos ne peut être
consideré comme point d'appui / au sens d'Euler »). « aux » est écrit
sur un autre mot ; « couples » sur « couple ». Dans le premier couple
d'équations, « $\frac{dz}{dx}$ » est écrit sans « $=0$ ». « indique »
est barré, « caractorise » écrit à la suite sur la ligne. « P. 157 »
est barré et « N.o 47 » écrit au-dessus ; la lecture « 157 » n'est pas
sûre, le dernier chiffre étant pris dans la biffure. L'initiale de « A
l'égard » est repassée. Dans la dernière paire d'équations, le second
crochet est écrit « $(1-e^{\omega})$ » devant $\text{sin}\,u\omega$, là où
les deux équations dont elle est tirée portent « $(1+e^{\omega})$ » ; le
feuillet est suivi. Le « ou » est écrit à droite, au bout de la seconde
équation. Une tache d'encre sous cette équation. « Néanmoins, » est
écrit au-dessus de « Cependant », barré.}

\page{424}

On doit conclure delà que lorsqu'il y a un nombre pair de points
de repos sur la lame élastique dont les extrêmités sont libres en
vertu des équations $\frac{d^2z}{dx^2}=0$ et $\frac{d^3z}{dx^3}=0$; elle vibre comme deux
lames séparés de moitié moins longues qui auroient chacune
une de leurs extrêmités assujettie aux mêmes conditions que celles
delalame entière, tandis que leur autre extrêmité seroit libre
en vertu des conditions $\frac{dz}{dx}=0$ et $\frac{d^3z}{dx^3}=0$.

Lorsqu'il y a un nombre impair de points de repos, la
valeur de $\underline{z}$ est donnée par l'équation
\[ z = \ldots\, e^{u\omega} - e^{(1-u)\omega} + (1+e^{\omega})\,\text{sin}.u\omega + (1-e^{\omega})\,\text{cos}\,u\omega \]
et d'après ceque j'ai dit pour le cas précédent, on voit
qu'un point de repos ne peut être point d'appui au
sens d'Euler que lorsqu'il satisfait aux deux conditions
\[ \left.\begin{array}{l} e^{u\omega} - e^{(1-u)\omega} + (1+e^{\omega})\,\text{sin}.u\omega + (1-e^{\omega})\,\text{cos}\,u\omega = 0 \\ e^{u\omega} - e^{(1-u)\omega} - \{(1+e^{\omega})\,\text{sin}.u\omega + (1-e^{\omega})\,\text{cos}\,u\omega\} = 0 \end{array}\right\} \text{ ou } \]
\[ \begin{array}{l} e^{u\omega} - e^{(1-u)\omega} = 0 \\ (1+e^{\omega})\,\text{sin}.u\omega + (1-e^{\omega})\,\text{cos}\,u\omega = 0 \end{array} \]
c'est adire que lorsque $u = \frac{1}{2}$.

On doit donc conclure que, lorsqu'il y a un nombre impair
de points de repos sur la lame élastique dont les extrêmités
sont libres au sens d'Euler, elle vibre comme deux lames
séparés, de moitié moins longues, qui auroient chacune
une de leurs extrêmités libres, tandis que leur autre extrêmité
seroit appuyée. Et en effet la condition $\frac{e^{\omega}-1}{e^{\omega}+1} = \text{tang}\,\frac{1}{2}\omega$,
qui resulte pour ce cas, de l'équation $\frac{2}{e^{\omega}+e^{-\omega}} = \text{cos}\,\omega$, laquelle
sert a déterminer l'angle $\omega$ lorsque les deux extrêmités sont
\note{En tête, au milieu, « (21) » ; en haut à droite, à l'encre, « 208 ».
Le dernier mot de la première ligne, « points », et le premier de la
seconde, « de repos », sont écrits sur d'autres mots, empâtés ; ce
qu'ils recouvrent ne se lit pas. « ceque » en un mot.
L'accolade qui réunit les
deux premières équations et le « ou » qui les relie aux deux suivantes
sont d'elle. La phrase se poursuit à la vue suivante.}

\page{425}

libres, n'est autre chose, en mettant $\underline{\omega}$ aulieu de $\frac{1}{2}\omega$, que
l'équation qui détermine le même angle $\underline{\omega}$ dans le second
cas examiné par Euler, c'est adire lorsqu'une extrêmité
est libre et l'autre appuyée. Il en est demême de
l'expression de $\underline{z}$, desorte qu'il est visible que ce grand
géomètre auroit pû se dispenser de calculer àpart
le cas dont il s'agit, puisque ce cas est expressement
contenu dans le premier

Dans le troisième cas d'Euler qui est celui d'une
extrêmité fixe et l'autre libre, on peut s'assurer,
en employant des raisonnemens semblables à ceux
dont j'ai fait usage pour le premier cas, que de
quelque manière que la lame se meuve, il ne
peut exister aucun autre point qui satisfasse aux
conditions des limites que les extrêmités reelles et
physiques de cette lame. Ainsi elle ne peut jamais
etre considerée \struck{y} comme separée en plusieurs
autres lames vibrant separément

Le quatrième cas d'Euler est celui des deux
extrêmités appuyés. Il est fort remarquable
en ce que, suivant \struck{la \uncertain{remarque}} l'observation de
cet illustre auteur, la lame se ploye alors
aux mêmes courbures que la corde tendue
\note{En tête, au milieu, « (22) » ; aucun nombre en haut à droite.
La première ligne achève la phrase de la vue 424 (« lorsque les deux
extrêmités sont / libres »). Le point manque après « le premier ». Le
mot biffé après « considerée » est une seule lettre, lue « y ». Après
« suivant », « la » est repassé et le mot qui suit est annulé par des
boucles serrées ; « remarque » est une lecture du sens plus que de la
lettre. Le point manque aussi après « separément ». La phrase se poursuit
à la vue suivante.}

\page{426}

Aussi est-il visible que, dans ce cas, tous les points de repos
sont de véritables limites, puisqu'ils satisfont \add{tous} aux conditions
$z=0$ et $\frac{ddz}{dx^2}=0$, et que, de plus, \struck{si} s'il y a, par exemple
$n-1$ points de repos, on aura \add{outre} les $n+1$ points \uncertain{d'extrêmités}
resultans des conditions $z=0$ et $\frac{ddz}{dx^2}=0$ $\underline{n}$ autres limites
analytiques assujetties aux conditions $\frac{dz}{dx}=0$ et $\frac{d^3z}{dx^3}=0$; desorte
que $2n+1$ exprimera le nombre des points qui pourront,
dans ce cas, etre considerés comme de veritables limites
des parties qu'ils separent.

Les deux derniers cas savoir, ceux d'une extrêmité fixe, l'autre
appuyée, et des deux extrêmités fixes, ont entr'eux le même
rapport que celui que j'ai remarqué entre les deux premiers
cas; et Euler auroit encore pû se dispenser de calculer
à part le v\textsuperscript{eme} qui se trouve compris dans le vi\textsuperscript{eme}.
A la vérité, la valeur de $\underline{z}$ qui appartient au v\textsuperscript{eme}
se présente sous une autre forme que celle qui seroit
déduite du vi\textsuperscript{eme}; mais cette différence vient uniquement
de ceque c'est l'extrêmité $x=0$ qu'Euler suppose
appuyée dans l'analyse de son v\textsuperscript{eme} cas, tandis que
dans le vi\textsuperscript{eme} ce ne peut jamais etre que le point
$x=\frac{1}{2}a$ qui soit censé appuyé. Au reste, dans ce
vi\textsuperscript{eme} cas, on trouve, comme dans le premier, qu'il
ne peut jamais y avoir qu'une seule extrêmité
\note{En tête, au milieu, « (23) » ; en haut à droite, à l'encre, « 209 ».
L'initiale de « Aussi » est repassée. « tous » est ajouté au-dessus de la
ligne, avec un signe d'insertion, après « satisfont ». « si » est barré
devant « s'il ». « outre » est écrit au-dessus de « les ». Le mot après
« points » est écrit sur un autre et se lit mal ; « d'extrêmités » est
offert avec doute. Le $n$ de « $\underline{n}$ autres limites » est
souligné. Les ordinaux « v\textsuperscript{eme} » et
« vi\textsuperscript{eme} » sont en chiffres romains minuscules, la
terminaison en exposant. La phrase se poursuit à la vue suivante.}

\page{427}

purement analytique, laquelle est toujours aumilieu de
la lame, et que lorsqu'il y a un nombre impair de points
de repos, cette extrêmité est appuyée en vertu des
conditions $z=0$ et $\frac{ddz}{dx^2}=0$, tandis que, lorsqu'il y a un
nombre pair de pareils points, l'extrêmité analytique
est libre en vertu des conditions $\frac{dz}{dx}=0$ et $\frac{d^3z}{dx^3}=0$.

D'après les remarques précédentes, il est maintenant
hors de doute que les points de repos des lames
élastiques ne doivent pas etre assimilés aux nœuds
de vibrations des cordes tendues; mais \struck{il est étonnant
que}, quoique la vraie théorie du mouvement de
ces lames ait été donnée depuis longtems par
Euler, que le cas des extrêmités libres ait été
le sujet d'un excellent memoire de Giordano
Ricati (Delle vibrazioni de cilindri T. 1 des M. dela
Société italienne) dans lequel cet auteur a calculé
avec beaucoup d'exactitude les differentes valeurs
de l'angle $x\omega$, et qu'en dernier lieu M\textsuperscript{r}
Chladni ait confirmé par l'experience les
résultats dela théorie qu'il reconnoit pour
entièrement satisfesante; \add{il est étonnant que} ni l'un ni l'autre
de ces savans auteurs ne se soient apperçu
que les intervalles qui separent les points
de repos sont inégaux entr'eux. \struck{il}
\note{En tête, au milieu, « (24) » ; aucun nombre en haut à droite. La
première ligne achève la phrase de la vue 426 (« qu'une seule
extrêmité / purement analytique »). « il est étonnant » est barré en
fin de ligne, et « que » au début de la suivante ; la même tournure,
« il est étonnant que », est ajoutée plus bas au-dessus de la ligne,
après « satisfesante », avec un signe d'insertion, où la phrase la
demande. La lettre devant $\omega$ dans « l'angle $x\omega$ » est lue $x$.
« Ricati », « de cilindri » : ainsi. Le dernier mot, « il », est
traversé d'un trait qui semble le biffer ; le feuillet suivant de la
même suite (vue 429, page (25)) s'ouvre sur « En effet », qui ne
continue pas « il ».}

\page{428}

\[ f^2 = dx^2 + \left(\frac{dz}{dx}\right)^2 dx^2 \]
\[ g^2 = dx^2 + \left\{\left(\frac{dz}{dx}\right)dx + \left(\frac{d^2z}{dx^2}\right)dx^2\right\}^2 \]
\[ h^2 = 4dx^2 + \left\{2\left(\frac{dz}{dx}\right)dx + \left(\frac{d^2z}{dx^2}\right)dx^2\right\}^2 \]
\[ f^2 + g^2 - h^2 = -2dx^2 - \uncertain{2}\,dx^2\left(\frac{dz}{dx}\right)\left(\frac{d^2z}{dx^2}\right)dx \]
\[ (f^2 + g^2 - h^2) = 4dx^4 + 8dx^5\left(\frac{dz}{dx}\right)^2\left(\frac{d^2z}{dx^2}\right) \]
\[ \qquad + 4dx^6\left(\frac{dz}{dx}\right)^4\left(\frac{d^2z}{dx^2}\right)^2 \]
\[ 4f^2g^2 = 4dx^4 \]

\[ {}^{d}\iint\left(\frac{d^2z}{dx^2} + \frac{d^2z}{dy^2}\right)\delta\frac{d^2z}{dx^2}\,dx\,dy \]
\[ = \int\ill{}\left(\frac{d^2z}{dx^2} + \frac{d^2z}{dx^2}\right)\delta\frac{dz}{dx}\,dy - \iint\frac{d\left(\frac{d^2z}{dx^2} + \frac{d^2z}{dy^2}\right)}{dx}\,\delta\frac{dz}{dx}\cdot dy \]
\[ \iint\frac{\left(\frac{d^2z}{dx^2} + \frac{d^2z}{dy^2}\right)\delta\ill{}}{dx} \]
\note{Petit feuillet d'un autre papier, plus court et plus étroit que les
feuillets qui l'entourent, monté sur onglet, le bord gauche arraché ; en
haut à droite, à l'encre, « 210 ». Au bas, le cachet « Bibliothèque
impériale MSS ». Des calculs sans un mot, d'une plume plus fine et d'une
écriture plus petite que la mise au net des vues 421–427 ; ils ne
continuent pas la phrase de la vue 427. Les accolades sont d'elle ; les
parenthèses droites de plusieurs facteurs sont à peine formées. Le
coefficient du second terme de $f^2+g^2-h^2$ est un chiffre bouclé,
lu 2 avec doute, suivi d'une tache d'encre. Le carré de
$(f^2+g^2-h^2)$ n'est pas écrit au premier membre, bien que le second en
soit le carré ; et les exposants 2 et 4 de $\left(\frac{dz}{dx}\right)$
aux deux derniers termes sont ceux du feuillet. Le « d » isolé devant la
première intégrale double est écrit plus haut que la ligne. Le signe
qui ouvre la deuxième ligne est un $\int$ suivi d'un trait appuyé
illisible, peut-être un second $\int$ ou un $\delta$ repassé. Dans la
parenthèse de cette ligne, le second terme est écrit
$\frac{d^2z}{dx^2}$, là où la première ligne porte $\frac{d^2z}{dy^2}$.
Le $\delta$ est le « d » à hampe bouclée de ses calculs de variation.
Sous le second terme, « $\int\frac{dz}{dx}\cdot dy$ » est récrit à
part, précédé d'un signe d'intégrale. La dernière expression est écrite
en partie sur le cachet ; le signe qui suit le $\delta$ du numérateur ne
se lit pas.}

\page{429}

En effet Giordano Ricati n'auroit eu qu'à comparer entr'elles
les valeurs qu'il a calculés \add{pour remarquer} qu'une partie comprise entre l'extrêmité
et le point de repos le plus voisin doit etre moins grande
que la moitié de l'espace compris entre le même nœud et
le suivant, que la partie comprise entre le second
nœud et le troisième doit etre moins grande que celle
qui est limitée par les deux premiers nœuds, et ainsi
de suite; desorte qu'en général une partie comprise
entre deux nœuds doit etre dautant plus étendue qu'elle
est plus éloignée de l'extrêmité et plus voisine du milieu
delalame quoique pourtant les differences soient
beaucoup plus grandes entre les parties qui sont près
des extrêmités qu'entre celles qui s'en éloignent davantage.
On ne peut douter que M\textsuperscript{r} Chladni n'eut
remarqué aussi les mêmes differences, \struck{si} s'il eut employé
des lames suffisamment longues; car je me suis
assuré que l'experience les confirme\struck{\ill{}} avec
la plus grande exactitude.

Euler lui même n'a sans doute pas fait la remarque
dont je viens de parler; car en voulant donner N.o 47
de son memoire un exemple des effets de l'application
d'un stilet dans un point quelconque d'une lame
dont les extrêmités seroient libres, fixés ou appuyés, il
\note{En tête, au milieu, « (25) », barré d'un trait oblique ; en haut à
droite, à l'encre, « 211 ». Dans la marge de gauche, le cachet
« Bibliothèque impériale MSS ». « pour remarquer » est écrit au-dessus
de la ligne, entre « calculés » et « qu'une ». La fin de « confirme »
est surchargée et barrée ; les lettres biffées ne se lisent pas. La
phrase se poursuit à la vue suivante. La vue 431 photographie une
seconde fois ce même feuillet ; la vue 430 est le dos, blanc, de la
feuille grise sur laquelle est monté le petit feuillet de la vue 428.}

\page{432}

établit une équation entre les valeurs de $\frac{dz}{dx}$ qui conviennent
aux deux branches séparées par l'application qu'il suppose;
or, cette équation qui \struck{\ill{}} a lieu en effet pour le cas
des extrêmités appuyés qu'il a choisi pour exemple,
ne sauroit subsister lorsque les deux branches ont des
courbures differentes; desorte que toute cette analyse
qui d'ailleurs seroit sujette à quelqu'objection ne
pourroit s'appliquer en general, aux autres états
des extrêmités. Au reste on jugera peut-etre cette
recherche surabondante si on observe que
l'effet de l'application d'un obstacle \struck{doit}
qui laisse subsister la communication du mouvement
entre les deux portions delalame, doit toujours
se borner à déterminer une des manières de
vibrer dont elle est susceptible à l'exclusion
des autres, qu'un pareil obstacle peut etre placé
dans un point de repos quelconque, et qu'il arrive
même toujours, lorsque les extrêmités sont libres
ou fixes, que ce n'est pas dans un véritable
point d'appui au sens d'Euler qu'il convient
de serrer la lame pour obtenir les
differens sons dont elle est susceptible.
\note{En tête, au milieu, « (26) » ; aucun nombre en haut à droite. La
première ligne achève la phrase de la vue 429 (« il / établit une
équation »), par-dessus la vue 430, blanche, et la vue 431, qui répète la
vue 429. Le mot barré après « qui », à la troisième ligne, est court et
noirci ; il ne se lit pas. « doit » est barré en fin de ligne après
« obstacle ». À la dernière ligne, « elle est » est repassé d'une encre
plus épaisse. L'initiale de « dans un point de repos » est tachée.}

\page{433}

D'après ceque je viens de dire, on voit que rien n'est plus vague
que la signification des mots \emph{libre} et \emph{appuyé} appliqués aux
differens points de la lame élastique; car après avoir \struck{\ill{}} remarqué
qu'il existe deux manières d'etre analytiquement differentes qui
peuvent appartenir aux extrêmités libres, je trouve encore que
l'appui physique est souvent placé dans un point qui ne
peut etre consideré comme appuyé au sens d'Euler

N.o 9 Je passe enfin à l'examen des conditions auxquelles sont
assujettis les points du contour des plaques élastiques vibrantes.
Si dans la formule N.o 1 résultante de la differentiation par
parties, on prend les termes qui ne sont affectés que d'un
seul signe d'intégration, que l'on y mette pour \uncertain{$E$.} sa
valeur, et que l'on divise par $\uncertain{\beta}^2$ et par les facteurs
dépendans de l'épaisseur et dela nature de la matière,
quantités \struck{qui} dont aucune ne peut etre supposée nulle,
on trouvera que les conditions des limites se reduisent aux
deux équations $S\left\{\left(\frac{ddz}{dx^2}\right) + \left(\frac{ddz}{dy^2}\right)\right\}\delta.\frac{dz}{dx\,dy} = 0$, $S\frac{d\left\{\left(\frac{ddz}{dx^2}\right) + \left(\frac{ddz}{dy^2}\right)\right\}}{dx\,dy}\,\delta z = 0$.

Afin de démêler plus aisément ceque signifient ces équations,
j'examinerai les conséquences que l'on peut en tirer dans les
hypothèses les plus simples. Je supposerai donc d'abord que
la surface vibrante soit rectangulaire, que le plan de cette
surface en repos soit pris pour celui des $\underline{x}$ et $y$ afin que
l'origine soit en un des points de repos dela même surface
\note{En tête, au milieu, « (27) », barré d'un trait oblique comme
« (25) » à la vue 429 ; en haut à droite, à l'encre, « 212 ». « libre »
et « appuyé » sont soulignés. Le mot barré devant « remarqué » est un
seul signe empâté. Le point manque après « au sens d'Euler ». « N.o 9 »
est écrit en tête de l'alinéa, en retrait, au-dessus de la ligne : un
numéro de paragraphe, le premier de ce lot ; « N.o 1 », dans la ligne,
renvoie à une formule d'une page qui n'est pas dans ce lot. La lettre
après « pour » est une capitale suivie d'un point, lue « E » avec doute ;
la lettre grecque sous le carré, lue $\beta$, n'est pas sûre non plus.
Le signe « S » qui ouvre les deux équations est son signe de somme
intégrale, gardé tel quel. Dans la première, « $dx\,dy$ » est écrit sous
la barre de $\frac{dz}{\ }$, là où la vue 428 écrit
$\delta\frac{dz}{dx}\,dy$. Le texte se poursuit à la vue suivante.}

\page{434}

enfin que chacun des côtés soit parallelle à un des axes \add{des} $\underline{x}$ et $y$.
cela posé, la suite des points qui servent de limites aux
valeurs de $\underline{x}$, c'est-à-dire ceux qui appartiennent aux côtés
parallelles à l'axe des $y$. devront satisfaire aux conditions
\[ S\frac{1}{dy}\left\{\left(\frac{ddz}{dx^2}\right) + \left(\frac{ddz}{dy^2}\right)\right\}\delta\left(\frac{dz}{dx}\right) = 0, \quad S\frac{1}{dy}\,\frac{d\left\{\left(\frac{ddz}{dx^2}\right) + \left(\frac{ddz}{dy^2}\right)\right\}}{dx}.\delta z = 0, \]
en vertu de l'un ou l'autre des quatre couples d'équations:
\[ z=0 \text{ et } \left(\frac{dz}{dx}\right)=0; \quad z=0 \text{ et } \left(\frac{ddz}{dx^2}\right) + \left(\frac{ddz}{dy^2}\right) = 0; \]
\[ \frac{d\left\{\left(\frac{ddz}{dx^2}\right) + \left(\frac{ddz}{dy^2}\right)\right\}}{dx} = 0 \text{ et } \left(\frac{dz}{dx}\right)=0; \quad \frac{d\left\{\left(\frac{ddz}{dx^2}\right) + \left(\frac{ddz}{dy^2}\right)\right\}}{dx} = 0 \text{ et } \left(\frac{ddz}{dx^2}\right) + \left(\frac{ddz}{dy^2}\right) = 0 \]
l'analogie est frapante entre ceque l'on trouve ici et
cequi a lieu pour les extrêmités de la simple lame.
On voit en effet que les deux derniers \struck{groupes} \add{couples} ne peuvent
appartenir qu'aux points libres du contour, tandis que
les deux premiers conviennent aux points de repos du
même contour. cependant on ne peut plus regarder
le point auquel \uncertain{appartiendroient} les conditions $z=0$ et $\left(\frac{dz}{dx}\right)=0$
comme fixe parcequ'elles ne suffisent pas pour que le
plan tangent tombe dans le plan des $\underline{x}$ et $y$ et qu'il
faudroit encore pour cela que l'on eut $\left(\frac{dz}{dy}\right)=0$.
On voit demême que la suite des points qui servent
de limites aux valeurs de $y$, c'est adire ceux qui appartiennent
aux côtés parallelles à l'axe des $\underline{x}$ doivent satisfaire
aux conditions $S\frac{1}{dx}\left\{\left(\frac{ddz}{dx^2}\right) + \left(\frac{ddz}{dy^2}\right)\right\}\delta.\left(\frac{dz}{dy}\right) = 0$, $S\frac{1}{dx}\,\frac{d\left\{\left(\frac{ddz}{dx^2}\right) + \left(\frac{ddz}{dy^2}\right)\right\}}{dy}\,\delta z = 0$.
en vertu de l'un ou l'autre des quatre couples d'équations:
\note{En tête, au milieu, « (28) », non barré ; aucun nombre en haut à
droite. La première ligne continue la phrase de la vue 433. « des » est
ajouté au-dessus de la ligne après « axes », avec un signe d'insertion. « groupes » est barré,
« couples » écrit au-dessus. Le verbe après « auquel » est surchargé à
sa fin : « appartiendroient » est offert avec doute. Les initiales de
« comme » et de « On voit demême » sont empâtées. Le second couple de la
page s'arrête sur les deux-points ; les quatre couples pour $y$ sont à la
vue suivante.}

\page{435}

\[ z=0 \text{ et } \left(\frac{dz}{dy}\right)=0; \quad z=0 \text{ et } \left(\frac{ddz}{dx^2}\right) + \left(\frac{ddz}{dy^2}\right) = 0; \quad \frac{d\left\{\left(\frac{ddz}{dx^2}\right) + \left(\frac{ddz}{dy^2}\right)\right\}}{dy} = 0 \]
\[ \text{et } \left(\frac{dz}{dy}\right)=0; \quad \frac{d\left\{\left(\frac{ddz}{dx^2}\right) + \left(\frac{ddz}{dy^2}\right)\right\}}{dy} = 0 \text{ et } \left(\frac{ddz}{dx^2}\right) + \left(\frac{ddz}{dy^2}\right) = 0. \]
Puisqu'il s'agit ici des surfaces rectangulaires, on voit que
les points situés aux quatre angles sont assujettis à quatre
conditions: \struck{x} car ils appartiennent également aux côtés parallelles
à l'axe des $\underline{x}$ et aux côtés parallelles à l'axe des $y$.
je \struck{les} nommerai points de \emph{double limite} tout point
qui remplira àla fois les conditions des limites des valeurs
de $\underline{x}$ et celles des valeurs de $y$. Ainsi les équations
$z=0$ et $\left(\frac{ddz}{dx^2}\right) + \left(\frac{ddz}{dy^2}\right) = 0$ qui forment le caractère analogue
à celui qu'Euler attribue aux points d'appui, ne pourront
appartenir qu'aux points de \emph{doubles limites}. \struck{si dans un
point} On aura encore un point de double limite, si on
peut satisfaire àla fois aux deux couples d'équations
$z=0$ et $\left(\frac{dz}{dx}\right)=0$; $z=0$ et $\left(\frac{dz}{dy}\right)$, et le plan tangent a
ce point tombera dans le plan de $\underline{x}$ et $y$.
au contraire le plan tangent sera seulement parallelle au
même plan des $\underline{x}$ et $y$ lorsque les conditions du
point de \emph{double limite} seront les deux couples d'équations
\[ \frac{d\left\{\left(\frac{ddz}{dx^2}\right) + \left(\frac{ddz}{dy^2}\right)\right\}}{dy} = 0 \text{ et } \left(\frac{dz}{dy}\right) = 0; \quad \frac{d\left\{\left(\frac{ddz}{dx^2}\right) + \left(\frac{ddz}{dy^2}\right)\right\}}{dx} = 0 \text{ et } \left(\frac{dz}{dx}\right) = 0. \]
Il est clair que tout point qui satisfera aux conditions
des limites, appartiendra à une ligne de limite analytique:
mais deplus on peut observer que si une pareille ligne
\note{En tête, au milieu, « (29) », barré d'un trait oblique ; en haut à
droite, à l'encre, « 213 ». Les deux premières lignes sont les quatre
couples d'équations annoncés à la fin de la vue 434, pour les côtés
parallèles à l'axe des $x$. Une petite croix, barrée, précède « car »
après « conditions: » ; aucun renvoi ne lui répond sur la page.
« les » est biffé après « je ». « double limite », « doubles limites »
sont soulignés. « si dans un point » est barré à la fin d'une ligne et
au début de la suivante ; l'initiale de « On » est repassée. Dans le
second couple, « $\left(\frac{dz}{dy}\right)$ » est écrit sans « $=0$ ».
La phrase se poursuit à la vue suivante.}

\page{436}

de limite satisfait, par exemple, aux conditions
\[ S\frac{1}{dy}\left\{\left(\frac{ddz}{dx^2}\right) + \left(\frac{ddz}{dy^2}\right)\right\}\delta\left(\frac{dz}{dx}\right) = 0, \quad S\frac{1}{dy}\,\frac{d\left\{\left(\frac{ddz}{dx^2}\right) + \left(\frac{ddz}{dy^2}\right)\right\}}{dx}.\delta z = 0 \]
et que pour un ou plusieurs de ses points on
ait $dy=0$; il faudra que les quatre conditions
\[ \left(\frac{ddz}{dx^2}\right) + \left(\frac{ddz}{dy^2}\right) = 0, \quad \left(\frac{dz}{dx}\right) = 0, \quad \frac{d\left\{\left(\frac{ddz}{dx^2}\right) + \left(\frac{ddz}{dy^2}\right)\right\}}{dx} = 0, \text{ et } z=0. \]
soient remplis àla fois pour les mêmes points:
ce qui donnera le caractere des points \emph{d'origine};
desorte que l'on pourra prendre indifferemment
sur les plaques rectangulaires tout point qui
satisfera à ces quatre conditions pour l'origine
des coordonnées dela surface.

Lorsque l'on veut appliquer les conditions des limites
à d'autres figures que les rectangles, il doit etre souvent
plus commode de les transformer suivant les cas;
par exemple, si on veut considérer une plaque ronde,
on peut prendre $x = r\,\text{cos}\,\varphi$, $y = r\,\text{sin}\,\varphi$; et on trouvera,
en désignant par $z'$ et $z''$ les valeurs de $\underline{z}$ qui
conviendroient à chacune des équations du second degré
auxquelles peut s'abaisser celle de la surface élastique
et observant qu'aux limites le rayon $r$ étant constant
on peut \struck{\uncertain{$r=a$}} faire $r=a$:
\[ z = z' + z''. \quad \left(\frac{dz}{dx}\right) = -\frac{\text{sin}\,\varphi}{a}\left\{\left(\frac{dz'}{d\varphi}\right) + \left(\frac{dz''}{d\varphi}\right)\right\}, \quad \left(\frac{ddz}{dx^2}\right) + \left(\frac{ddz}{dy^2}\right) = \ldots\{z' - z''\} \]
\[ \frac{d\left\{\left(\frac{ddz}{dx^2}\right) + \left(\frac{ddz}{dy^2}\right)\right\}}{dx} = \ldots\,\frac{\text{sin}\,\varphi}{a}\left\{\left(\frac{dz''}{d\varphi}\right) - \left(\frac{dz'}{d\varphi}\right)\right\} \]
\note{En tête, au milieu, « (30) », non barré ; aucun nombre en haut à
droite. La première ligne continue la phrase de la vue 435 (« si une
pareille ligne / de limite satisfait »). Après « ait $dy=0$ », un signe
entre virgule et point-virgule. « d'origine » est souligné. « sin » est
écrit avec un s long, « cos » et « sin » suivis d'un point à la
deuxième formule. Le groupe barré après « on peut » est court et pris
dans une croix ; on croit y lire « $r=a$ », que la suite récrit. Les
points de suspension, « $=\ldots$ », sont d'elle : elle laisse en blanc
le facteur qui précède l'accolade, comme sur la copie d'Euler. La
page s'arrête sur la dernière équation ; elle ne porte pas de mot de
réclame.}

\page{437}

\[ \frac{d\left\{\left(\frac{ddz}{dx^2}\right) + \left(\frac{ddz}{dy^2}\right)\right\}}{dy} = \ldots\,\frac{\text{cos}\,\varphi}{a}\left\{\left(\frac{dz''}{d\varphi}\right) - \left(\frac{dz'}{d\varphi}\right)\right\} \]
Ainsi les conditions auxquelles sont assujettis les points qui servent
de limites aux valeurs de $\underline{x}$ seront
\[ S\frac{1}{dy}\left\{\left(\frac{ddz}{dx^2}\right) + \left(\frac{ddz}{dy^2}\right)\right\}\delta\left(\frac{dz}{dx}\right) = \struck{\ldots}\,S\frac{1}{d\varphi}(z'-z'')\,\frac{\text{tang}\,\varphi}{a}\,\delta\left\{\left(\frac{dz'}{d\varphi}\right) + \left(\frac{dz''}{d\varphi}\right)\right\} = 0 \]
\[ S\frac{1}{dy}\,\frac{d\left\{\left(\frac{ddz}{dx^2}\right) + \left(\frac{ddz}{dy^2}\right)\right\}}{dx}.\delta z = \ldots - S\frac{1}{d\varphi}\,\frac{\text{tang}\,\varphi}{a}\left\{\left(\frac{dz''}{d\varphi}\right) - \left(\frac{dz'}{d\varphi}\right)\right\}\delta(z'+z'') = 0. \]
Les mêmes substitutions donneront pour les conditions auxquelles
sont assujettis les points qui servent de limite aux valeurs de $y$
\[ S\frac{1}{dx}\left\{\left(\frac{ddz}{dx^2}\right) + \left(\frac{ddz}{dy^2}\right)\right\}\delta\left(\frac{dz}{dy}\right) = \ldots\,S\frac{1}{d\varphi}(z'-z'')\,\frac{\text{cot}.\varphi}{a}\,\delta\left\{\left(\frac{dz'}{d\varphi}\right) + \left(\frac{dz''}{d\varphi}\right)\right\} = 0 \]
\[ S\frac{1}{dx}\,\frac{d\left\{\left(\frac{ddz}{dx^2}\right) + \left(\frac{ddz}{dy^2}\right)\right\}}{dy}.\delta z = \ldots\,S\frac{1}{d\varphi}\struck{(z'-z'')}\,\frac{\text{cot}.\varphi}{a}\left\{\left(\frac{dz'}{d\varphi}\right) - \left(\frac{dz''}{d\varphi}\right)\right\}\delta(z'+z'') = 0. \]
Les premières seront donc remplies par la supposition
$\text{tang}\,\varphi = 0$, et les secondes par la supposition équivalente
$\text{cot}.\varphi = 0$., c'est adire que ces points seront aux extrêmités
des axes des $\underline{x}$ et des $y$, et parconséquent à angles droits
les uns des autres: c'est ceque l'on auroit pû conclure immed\add{iatement}
dela figure de cette plaque, \struck{laquelle} ne permet \struck{jamais} pas
qu'un seul et même point remplisse àla fois les deux
systemes de conditions qui caractérisent ceque j'ai nommé
point de \emph{double limite}

Pour tout autre point du contour que ceux dont je viens
de parler, les conditions des limites se reduisent à
\[ S(z'-z'')\,\delta\left\{\left(\frac{dz'}{d\varphi}\right) + \left(\frac{dz''}{d\varphi}\right)\right\} = 0 \qquad S\left\{\left(\frac{dz''}{d\varphi}\right) - \left(\frac{dz'}{d\varphi}\right)\right\}\delta(z'+z'') = 0, \]
et on peut y satisfaire par l'un ou l'autre des couples
\note{En tête, au milieu, « (31) », barré ; en haut à droite, à l'encre,
« 214 ». La première équation continue la liste de la vue 436. Dans la
première des quatre équations suivantes, un petit groupe après « $=$ »
est biffé, apparemment les points de suspension qu'elle écrit
ailleurs ; dans la dernière, « $(z'-z'')$ » est biffé devant
$\frac{\text{cot}.\varphi}{a}$. Le second terme de cette dernière est
écrit $\left(\frac{dz'}{d\varphi}\right) - \left(\frac{dz''}{d\varphi}\right)$,
dans l'ordre inverse de celui de la deuxième ; le feuillet est suivi.
« immed » est coupé au bord droit sans point d'abréviation ; la fin du
mot, « iatement », est ajoutée ici entre crochets éditoriaux. « laquelle »
et « jamais » sont barrés d'un trait ; la phrase reste sans relatif. Le
point manque après « double limite », souligné. La phrase se poursuit à
la vue suivante.}

\page{438}

d'équations $z'-z''=0$ et $z'+z''=0$; $z'-z''=0$ et $\left(\frac{dz'}{d\varphi}\right) - \left(\frac{dz''}{d\varphi}\right) = 0$;
$z'+z''=0$ et $\left(\frac{dz'}{d\varphi}\right) + \left(\frac{dz''}{d\varphi}\right) = 0$; $\left(\frac{dz'}{d\varphi}\right) - \left(\frac{dz''}{d\varphi}\right) = 0$ et $\left(\frac{dz'}{d\varphi}\right) + \left(\frac{dz''}{d\varphi}\right) = 0$.

Je ne m'étendrai pas davantage sur ce sujet: mais je
ne puis m'empêcher d'observer que la considération
des conditions auxquelles les points du contour doivent
satisfaire, est dela plus grande importance pour déterminer
le choix des intégrales particulières qui peuvent convenir
dans chaque cas, et qu'elle fournit aussi le moyen
de renfermer dans les mêmes formules l'analyse de
divers cas qui sembleroient, au premier coup d'œil, devoir
etre calculés à part. En effet, d'après ceque j'ai dit
plus haut, on voit que la suite des points qui satisfont
aux conditions des \struck{extrêmités} limites, forment sur la surface, des
lignes de limites analytiques, et que l'on peut
parconséquent concevoir ces surfaces comme reellement séparées
en autant de surfaces partielles dont les points de
contour peuvent etre dans un état different de ceux
de la plaque entière, sans que \add{pour cela} l'équation qui
convient au mouvement de cette plaque entière,
cesse d'être applicable aux portions dans lesquelles
elle est analytiquement separée.
\note{En tête, au milieu, « (32) », non barré ; aucun nombre en haut à
droite. Les deux premières lignes achèvent la phrase de la vue 437
(« par l'un ou l'autre des couples / d'équations »). Après « des », un
mot commencé, lu « extrêmités », est surchargé par « limites » ; une
tache d'encre au-dessous. « pour cela » est écrit au-dessus de la ligne,
après « sans que ». Le texte s'arrête ici, aux deux tiers de la page,
sur un point : la fin d'une section. Le feuillet suivant (vue 439) ouvre
une section nouvelle.}

\page{439}

§. 4. Comparaison dela théorie à l'experience; \struck{\ill{}} \add{dans le cas des} \struck{cas} \add{des} plaques quarrées

N.o 10. En prenant d'abord la valeur de $\underline{z}$ donnée par l'équation
(C) N.o 3, et y introduisant les coëfficiens déterminés par Euler (P. 118 et suiv.
de son m.) pour le cas des extrêmités libres, en vertu des conditions
$\frac{ddz}{dx^2}=0$ et $\frac{d^3z}{dx^3}=0$ ou $\frac{ddz}{dy^2}=0$ et $\frac{d^3z}{dy^3}=0$, on trouvera
la formule
\[ z = \ldots\left\{\mathcal{A}\left\{e^{\frac{\omega x}{A}} \pm e^{\frac{(A-x)\omega}{A}} + (1\mp e^{\omega})\,\text{sin}.\frac{\omega x}{A} + (1\pm e^{\omega})\,\text{cos}.\frac{\omega x}{A}\right\}\right. \]
\[ \left. + B\left\{e^{\frac{\omega y}{A}} \pm e^{\frac{(A-y)\omega}{A}} + (1\mp e^{\omega})\,\text{sin}.\frac{\omega y}{A} + (1\pm e^{\omega})\,\text{cos}.\frac{\omega y}{A}\right\}\right. \]
dans laquelle les signes superieurs appartiennent au cas où il y a
un nombre pair de points de repos sur la simple lame, et \add{où} les signes
inferieurs conviennent à celui où il y a un nombre impair de pareils
points sur la même lame.

Dans l'un et l'autre cas il est visible qu'en fesant $\mathcal{A}$ ou
$B=0$, il devra y avoir autant \add{de lignes} nodales droites sur la plaque
quarrée qu'il y auroit \struck{\ill{}} de points de repos sur la simple lame
dont la plaque ne diffère alors que par l'étendue d'une
dimension qui n'entre plus dans le calcul.

Il \struck{est \ill{}} faut àprésent distinguer le cas des signes superieurs
de celui des signes inferieurs. Pour cela, en \struck{supp}\uncertain{fesant} d'abord
$\mathcal{A} = -B$. j'écris
\[ \left.\begin{array}{l} z = \ldots\,\mathcal{A}\left\{e^{\frac{\omega x}{A}} + e^{\frac{(A-x)\omega}{A}} + (1-e^{\omega})\,\text{sin}.\frac{\omega x}{A} + (1+e^{\omega})\,\text{cos}.\frac{\omega x}{A}\right\} \\ \qquad - \left\{e^{\frac{\omega y}{A}} + e^{\frac{(A-y)\omega}{A}} + (1-e^{\omega})\,\text{sin}.\frac{\omega y}{A} + (1+e^{\omega})\,\text{cos}.\frac{\omega y}{A}\right\} \end{array}\right\}\ldots(C') \]
On voit que le plan des $\underline{x}$ et $y$ sera alors celui dela plaque en
repos; car $\underline{z}$ sera nulle toutes les fois que $\underline{x}$ égalera $y$; cequi
déterminera une première ligne nodale diagonale sur la surface
quarrée; $\underline{z}$ sera également nulle toutes les fois que
\note{En tête, au milieu, « (33) », barré ; en haut à droite, à l'encre,
« 215 ». Une section nouvelle commence : le titre, souligné d'un court
trait au milieu de la ligne suivante, ouvre le « §. 4 » — lu sur un
signe de paragraphe tracé comme « S. », suivi d'un point — et le
paragraphe est numéroté « N.o 10 », à la suite du « N.o 9 » de la
vue 433. Dans le titre, après le point-virgule, un mot est barré de
boucles serrées et ne se lit pas ; « dans le cas des » est écrit
au-dessus, puis « cas » est barré et « des » récrit au-dessus de la
ligne. « (C) N.o 3 » renvoie à une équation d'un paragraphe qui n'est pas
dans ce lot ; « P. 118 et suiv. de son m. » à la page du mémoire d'Euler.
Le premier coefficient est une capitale cursive bouclée, distincte du $A$
des dénominateurs ; elle est rendue $\mathcal{A}$, et le second
coefficient $B$. Dans la première formule, au dernier terme, un $B$
empâté est écrit sur « (1 » ; on lit « $(1\pm e^{\omega})$ », comme à la
ligne suivante. « où » est ajouté au-dessus de « les signes » ; « de
lignes » au-dessus de « nodales ». Après « auroit », quelques lettres
sont noircies et ne se lisent pas. « est » et le mot qui le suit, devant
« faut », sont barrés ; le second ne se lit pas. Devant « d'abord »,
« supp » est barré et la fin du mot récrite dessus ; « fesant » est
offert avec doute. Dans (C'), le facteur $\mathcal{A}$ n'est pas répété
devant la seconde accolade ; l'étiquette « (C') » est écrite à droite de
la grande accolade. La phrase se poursuit à la vue 440.}

\page{440}

l'équation $y = A - x$ aura lieu. car à cause dela condition
$\text{tang}\,\frac{1}{2}\omega = \frac{1-e^{\omega}}{1+e^{\omega}}$ qui appartient à ce cas, l'équation (C') se
réduit à
\[ z = \ldots\,\mathcal{A}\left\{(1-e^{\omega})\,2\,\text{cos}.\frac{\omega}{2}\,\text{sin}\frac{(A-2x)\omega}{2A} - (1+e^{\omega})\,2\,\text{sin}.\frac{\omega}{2}\,\text{sin}\frac{(A-2x)\omega}{2A}\right\} = 0. \]
Il y aura donc une seconde ligne nodale diagonale sur
la plaque quarrée, lorsque la valeur de l'angle $\underline{\omega}$
indiquera un nombre pair de points de repos sur
la simple lame.

En continuant de faire $\mathcal{A} = -B$ mais prenant les signes
inferieurs qui conviennent au cas d'un nombre impair
de points de repos dans la simple lame, on trouve
qu'il n'y a alors qu'une ligne nodale diagonale, savoir celle
qui est formée de la suite des points pour lesquels $x = y$.

Les lignes nodales diagonales ne sont pas les seules
lignes nodales qui doivent exister sur les plaques
quarrées dont les mouvemens \struck{sont} appartiennent
à l'intégrale (C) dans la supposition $\mathcal{A} = -B$;
et il est visible que leur nombre et leur disposition
dependent des valeurs de l'angle $\underline{\omega}$.

En général, quelque soit le rapport entre les
valeurs de $\mathcal{A}$ et de $B$, on voit, à la simple
inspection de l'équation (C), que la figure
nodale doit passer par tous les points
d'intersection des lignes nodales qui conviendroient
aux suppositions successives $\mathcal{A} = 0$, $B = 0$; desorte
que, si on represente par $\underline{n}$ le nombre de
\note{En tête, au milieu, « (34) », non barré ; aucun nombre en haut à
droite. La première ligne achève la phrase de la vue 439 (« toutes les
fois que / l'équation $y = A - x$ aura lieu »). Dans
$\frac{1-e^{\omega}}{1+e^{\omega}}$, le $e$ du numérateur est à peine
formé. Le renvoi « (C') » de la deuxième ligne a son accent ; ceux de
« l'intégrale (C) » et de « l'équation (C) », plus bas, n'en ont pas, et
sont gardés ainsi. L'initiale de « Il y aura » et celle de « En
général » sont repassées. « sont » est barré devant « appartiennent ».
La phrase se poursuit au-delà de ce lot, à la vue 441.}

\end{document}
